\section*{Hoofdstuk \ref{ch:g9:genes-and-dna} --- Genen en DNA}

\begin{solution}{exo:g9:genes-and-dna:1}
Een \omterm{def:g9:genes-and-dna:gene}{gen} is een stuk van een \omterm{def:g9:chromosomes:chromosomes}{chromosoom} dat de instructies voor
één \omterm{def:g9:heredity-and-traits:traits}{kenmerk} draagt, op een vast adres; \omterm{def:g9:genes-and-dna:gene}{allelen} zijn de
uitvoeringen ervan --- A, B en O van het bloedgroepgen.
\end{solution}

\begin{solution}{exo:g9:genes-and-dna:2}
Omdat chromosomen in paren komen: van elk \omterm{def:g9:genes-and-dna:gene}{gen} rijdt op elke
partner één exemplaar mee --- één van de moeder, één van de
vader, op hetzelfde adres.
\end{solution}

\begin{solution}{exo:g9:genes-and-dna:3}
De volgorde van de chemische letters langs de streng. Vier
\omterm{def:g5:biodiversity:species}{soorten} letters, ongeveer drie miljard ervan in één menselijk
spel.
\end{solution}

\begin{solution}{exo:g9:genes-and-dna:4}
De strengen zijn letter voor letter elkaars aanvulling: uit
elkaar getrokken bouwt elke streng haar partner opnieuw op, en
van één ladder worden twee gelijke ladders --- een vol exemplaar
per dochtercel.
\end{solution}

\begin{solution}{exo:g9:genes-and-dna:5}
Een zeldzame kopieerfout --- veranderde letters die een nieuw
\omterm{def:g9:genes-and-dna:gene}{allel} maken. Meestal neutraal, soms schadelijk, af en toe
nuttig.
\end{solution}

\begin{solution}{exo:g9:genes-and-dna:6}
Alle \omterm{def:g6:exploring-our-environment:components}{levende wezens} schrijven en lezen hun \omterm{def:g9:genes-and-dna:gene}{genen} in dezelfde
DNA-letters en dezelfde code. Gevolgen: het sterkste bewijs voor
een gemeenschappelijke oorsprong --- en het feit dat \omterm{def:g9:genes-and-dna:gene}{genen}
tussen \omterm{def:g5:biodiversity:species}{soorten} verplaatsen werkt, omdat de taal overal past.
\end{solution}

\begin{solution}{exo:g9:genes-and-dna:7}
Elke ouder draagt bruin met blauw --- bruin komt tot uiting,
blauw zwijgt. Elke ouder deelt blauw uit met één kans op twee;
de uitdelingen zijn onafhankelijk, dus komt blauw-met-blauw met
één kans op vier terecht --- het blauwogige kind.
\end{solution}

\begin{solution}{exo:g9:genes-and-dna:8}
Groep A: A met A, of A met O. Groep B: B met B, of B met O.
Groep AB: A met B --- beide komen tot uiting. Groep O: O met
O.
\end{solution}

\begin{solution}{exo:g9:genes-and-dna:9}
Beide bevatten de hele bibliotheek --- dezelfde 46 draden. De
\omterm{def:g8:nervous-communication:neuron}{zenuwcel} leest de \omterm{def:g9:genes-and-dna:gene}{genen} van het zenuwapparaat en laat de
keratinehoofdstukken van de \omterm{def:g1:the-five-senses:sense}{huid} dicht; de huidcel doet het
omgekeerde. Specialisatie is een uitleenbeleid, geen andere
bibliotheek.
\end{solution}

\begin{solution}{exo:g9:genes-and-dna:10}
Pureer het fruit (de cellen worden geopend); zout water en
afwasmiddel (de membranen --- de omhulsels van de \omterm{def:g5:first-look-at-cells:cell}{cel} en van de
kern --- worden opgelost, het \omterm{def:g9:genes-and-dna:dna}{DNA} komt vrij en blijft opgelost);
zeef (het gruis eruit); koude sterke drank erbovenop (het \omterm{def:g9:genes-and-dna:dna}{DNA}
lost op de grens weer uit): witachtige, opwindbare draden --- het
molecuul zelf.
\end{solution}

\begin{solution}{exo:g9:genes-and-dna:11}
Familieniveau: de kin keert terug langs de lijnen van de
voortplanting, één generatie lang ongetoond meegedragen.
Chromosoomniveau: haar bepaler rijdt mee op één draad, tot
\omterm{def:g8:sexual-reproduction:gametes}{geslachtscellen} gehalveerd en bij de \omterm{def:g5:flower-to-fruit:steps}{bevruchting} hersteld, aan
sommige kinderen uitgedeeld en aan andere niet. Molecuulniveau:
het is een \omterm{def:g9:genes-and-dna:gene}{allel} --- een spelling van een gezichtsvormend \omterm{def:g9:genes-and-dna:gene}{gen} ---
sinds de \omterm{prop:g9:genes-and-dna:explains}{mutatie} die het ooit schreef, de eeuw door getrouw
gekopieerd.
\end{solution}

\begin{solution}{exo:g9:genes-and-dna:12}
De \omterm{prop:g9:genes-and-dna:explains}{mutatie} schrijft: kopieerfouten maken nieuwe spellingen ---
molecuulniveau, dit hoofdstuk. Het schudden deelt uit: halveren
en \omterm{def:g5:flower-to-fruit:steps}{bevruchting} geven elk kind een verse combinatie ---
chromosoomniveau, \cref{ch:g9:chromosomes} (en
\cref{prop:g8:sexual-reproduction:shuffle}). De uiting toont: van
de twee uitgedeelde exemplaren, wat er getoond wordt en wat
zwijgend meerijdt --- de regel van twee exemplaren uit dit
hoofdstuk. Alleen de \omterm{prop:g9:genes-and-dna:explains}{mutatie} maakt werkelijk nieuwe spellingen;
de andere twee hercombineren en onthullen wat de \omterm{prop:g9:genes-and-dna:explains}{mutatie} ooit
schreef.
\end{solution}

\begin{solution}{pb:g9:genes-and-dna:1}
\textbf{1.} Een \omterm{prop:g9:genes-and-dna:explains}{mutatie}: één keer een kopieerfout in het
kleurstofgen, in de geslachtscellijn van een voorouder ---
letters veranderd, de instructies bedorven, en de verkeerde
spelling sindsdien getrouw gekopieerd.

\textbf{2.} Omdat de werkende instructies van P de kleurstof
bouwen, wat de zwijgende partner ook is: p rijdt
\emph{ongetoond meegedragen} mee --- het woord uit het album, nu
moleculair.

\textbf{3.} Met p op beide partners bestaan er op geen van beide
adressen werkende instructies: het kleurstofapparaat --- het
product van het \omterm{def:g9:genes-and-dna:gene}{gen} --- wordt eenvoudig nooit gebouwd. Het
\omterm{def:g9:heredity-and-traits:traits}{kenmerk} is het ontbreken van een werkende spelling.

\textbf{4.} Dat hun kleurstofgenen herkenbaar hetzelfde \omterm{def:g9:genes-and-dna:gene}{gen} zijn
--- vergelijkbare volgordes in dezelfde code --- verwante kopieën
van één voorouderlijke tekst, in de eigen geschiedenis van elke
\omterm{def:g5:biodiversity:species}{soort} onafhankelijk kapotgegaan.

\textbf{5.} De vier uitdelingen: P-met-P, P-met-p, p-met-P,
p-met-p --- even waarschijnlijk. De kleurloze opbouw: één kans op
vier.

\textbf{6.} Het overslaan: twee dragers die kleurstof toonden
gaven door wat ze niet toonden, en de verborgen spellingen
troffen elkaar --- \cref{ex:g9:heredity-and-traits:skip}, met het
volledige mechanisme uitgevoerd.

\textbf{7.} Allemaal dragen ze één p --- de hele bijdrage van hun
moeder op dat adres is p. Geen van hen toont de aandoening: de P
van de partner bouwt bij elk van hen de kleurstof.

\textbf{8.} De p van oma reed mee op één \omterm{def:g9:chromosomes:chromosomes}{chromosoom} van een paar;
het halveren van haar \omterm{def:g8:sexual-reproduction:gametes}{geslachtscel} deelde die partner uit aan de
\omterm{def:g4:animal-reproduction:sexual}{eicel} of \omterm{def:g8:reproductive-systems:male}{zaadcel} die de dragende ouder maakte; het halveren van
die ouder deelde hem opnieuw uit aan de \omterm{def:g8:sexual-reproduction:gametes}{geslachtscel} die het
kleinkind maakte --- van draad naar \omterm{def:g8:sexual-reproduction:gametes}{geslachtscel} naar \omterm{prop:g8:sexual-reproduction:core}{zygote},
twee keer achtereen.

\textbf{9.} Familie: twee gewoon ogende ouders, een kind met één
kans op vier, dragerschap door de \omterm{def:g2:born-grow-die:cycle}{geboorte} bewezen. \omterm{def:g9:chromosomes:chromosomes}{Chromosoom}:
de partners van één paar die P en p dragen, gehalveerd en
opnieuw gecombineerd. Molecuul: één \omterm{def:g9:genes-and-dna:gene}{gen}, één werkend en één
verkeerd gespeld \omterm{def:g9:genes-and-dna:gene}{allel}, en het spellingsverschil beslist alles.

\textbf{10.} Voor elk oordeel: bij een drager is het kapot zijn
van p onzichtbaar --- de P van de partner dekt het af, dus
``meldt'' niets aan de opbouw van de drager de verkeerde
spelling. Wat niet gezien kan worden, kan niet weggerepareerd
worden; verborgen spellingen houden zich onbeperkt in stand.

\textbf{11.} Haar \omterm{def:g9:heredity-and-traits:traits}{kenmerk} (geen kleurstof) is geërfd ---
p-met-p --- maar het \emph{gevolg} ervan wordt geleefd in een
\omterm{def:g6:exploring-our-environment:components}{omgeving} vol zon: de ontbrekende kleurstof was de zonnebrandcrème
van het lichaam, dus schrijft de \omterm{def:g6:exploring-our-environment:components}{omgeving} meer op haar dan op
haar afgeschermde ouders --- de erfenis zet de marge, de wereld
vult haar in.

\textbf{12.} Wat er kopieert: de dubbele streng, waarvan elke
helft haar partner opnieuw opbouwt bij elke celdeling en elke
generatie. Wat er één keer verkeerd gekopieerd werd: een paar
letters van het kleurstofgen, in één oeroude \omterm{def:g8:sexual-reproduction:gametes}{geslachtscel}. Waarom
de fout blijft: het kopieerapparaat is getrouw --- het kopieert
wat er staat, verkeerde spellingen met dezelfde nauwkeurigheid
als zin.
\end{solution}
